\documentclass[11pt]{article}
\usepackage[margin=1in]{geometry}
\usepackage[T1]{fontenc}
\usepackage{lmodern}
\usepackage{amsmath,amssymb,amsthm}
\usepackage[hidelinks]{hyperref}
\usepackage{microtype}

\newtheorem*{theorem}{Theorem}

\title{A Salem Spectral Arc Answers Both Fourier-Dimension Questions\\
in arXiv:2002.03855}
\author{Literature-implied full-result note\\
Run \texttt{fa\_banach\_001}, lane 6}
\date{11 August 2026}

\begin{document}
\maketitle

\begin{center}
\fbox{\parbox{0.91\textwidth}{\textbf{Status: literature-implied full
answer.} Li and Liu's later Salem spectral arc directly answers the second
question.  A short restriction-entropy computation, apparently not made in
that paper, shows that the same measure also attains equality in Shi's
necessary dimension inequality and therefore answers the first question.}}
\end{center}

\section*{The two source questions}

At the end of Section~6 of Ruxi Shi,
\emph{On dimensions of frame spectral measures and their frame spectra},
arXiv:2002.03855 \cite{Shi}, the following questions are posed.

First, does there exist a singular continuous frame-spectral measure for
which equality holds in
\begin{equation}\label{eq:Shi}
 \sup_{\mu(K)>0,\,\mu(\partial K)=0}\dim_F\mu_K
 \ \leq\!
 \inf_{\mu(K)>0,\,\mu(\partial K)=0}
 \overline{\dim}_e\mu_K?
\end{equation}
Second, does there exist a singular continuous frame-spectral measure with
nonzero Fourier dimension?

Shi defines entropy dimension for probability measures.  Thus, as is
standard and as scale invariance requires, entropy dimension of a positive
restriction means that of the normalized restriction
\(\mu_K/\mu(K)\).  Fourier dimension is unchanged by multiplication by a
positive scalar.

\section*{The later arc}

Let \(I=[-1/2,1/2]\), let
\[
 \gamma(x)=(x,\sqrt{1-x^2}),
 \qquad \sigma=\gamma_*(dx|_I),
\]
so that
\[
 \int f\,d\sigma=\int_{-1/2}^{1/2}
 f(x,\sqrt{1-x^2})\,dx.
\]
Li and Liu exhibit exactly this measure in arXiv:2506.01280
\cite{LiLiu}.  They observe that it has orthonormal spectrum
\(\mathbb Z\times\{0\}\), and Section~10 proves
\begin{equation}\label{eq:decay}
 |\widehat\sigma(\xi)|\lesssim |\xi|^{-1/2}.
\end{equation}
Indeed, the pullback of
\(e^{-2\pi i (n,0)\cdot y}\) under \(\gamma\) is
\(e^{-2\pi i nx}\), and these functions form the standard orthonormal basis
of \(L^2(I)\).  Thus \(\sigma\) is spectral, hence frame-spectral.

The map \(\gamma\) is injective and Lipschitz on \(I\).  Consequently
\(\sigma\) has no atoms, while its support is a planar arc of zero planar
Lebesgue measure; hence it is singular continuous.  The arc has Hausdorff
dimension one.  Equation~\eqref{eq:decay} and the general inequality
\(\dim_F\sigma\leq\dim_H\operatorname{supp}\sigma\) give
\begin{equation}\label{eq:fourierone}
 \dim_F\sigma=1.
\end{equation}
This already answers Shi's second question affirmatively.

\section*{Upgrade: all positive restrictions have entropy dimension one}

\begin{theorem}
For the measure \(\sigma\) above,
\[
 \sup_{\sigma(K)>0,\,\sigma(\partial K)=0}\dim_F\sigma_K
 =1=
 \inf_{\sigma(K)>0,\,\sigma(\partial K)=0}
 \overline{\dim}_e\!\left(\frac{\sigma_K}{\sigma(K)}\right).
\]
In particular equality holds in \eqref{eq:Shi}.
\end{theorem}

\begin{proof}
Fix a Borel set \(K\subset\mathbb R^2\) with \(\sigma(K)>0\), and set
\[
 A=\gamma^{-1}(K)\subset I,\qquad a=|A|>0,
 \qquad \nu=\frac{\sigma_K}{\sigma(K)}.
\]
The first-coordinate marginal of \(\nu\) is
\(a^{-1}{\bf1}_A(x)\,dx\).  If \(J\) is a level-\(n\) dyadic interval,
then its marginal mass \(p_J\) satisfies
\[
 p_J=\frac{|A\cap J|}{a}\leq\frac{2^{-n}}a.
\]
For all sufficiently large \(n\), Shannon entropy therefore satisfies
\[
 H\bigl((\pi_1)_*\nu,\mathcal D_n^{(1)}\bigr)
 =\sum_J p_J\log_2\frac1{p_J}
 \geq \log_2\frac1{\max_Jp_J}
 \geq n+\log_2 a.
\]
The first-coordinate partition is a coarsening of the planar dyadic
partition, so
\begin{equation}\label{eq:lowerH}
 H(\nu,\mathcal D_n^{(2)})\geq n+O(1).
\end{equation}

Conversely, the graph of the Lipschitz function
\(x\mapsto\sqrt{1-x^2}\) on \(I\) meets at most \(C2^n\) level-\(n\)
dyadic squares.  A probability vector supported on at most \(N\) atoms has
entropy at most \(\log_2N\), whence
\begin{equation}\label{eq:upperH}
 H(\nu,\mathcal D_n^{(2)})\leq \log_2(C2^n)=n+O(1).
\end{equation}
Equations~\eqref{eq:lowerH}--\eqref{eq:upperH} show that every positive
restriction has entropy dimension exactly one.  The null-boundary condition,
needed in Shi's frame-restriction lemma, only narrows this class and does not
alter the conclusion.

Every \(\sigma_K\) is supported on the same one-dimensional Lipschitz arc,
so \(\dim_F\sigma_K\leq1\).  On the other hand
\(K=\mathbb R^2\) is admissible, has empty boundary, and gives
\(\sigma_K=\sigma\); hence \eqref{eq:fourierone} makes the supremum equal to
one.  The entropy calculation makes the infimum equal to one.
\end{proof}

\section*{Provenance and scope}

The Salem spectral arc, its orthonormal spectrum, and the decay estimate are
from Li--Liu.  Their paper cites Shi for a different frame-density argument,
but does not state that this arc answers Shi's questions.  The identification
with the second question and especially the all-restrictions entropy
calculation above are therefore recorded as a \emph{literature-implied}
answer, not as an author-explicit literature resolution and not as a new
independent construction.

The conclusion is planar, exactly as Shi's questions permit.  It does not
settle Li--Liu's later, stricter real-line question.

\begin{thebibliography}{9}
\bibitem{Shi}
R. Shi,
\emph{On dimensions of frame spectral measures and their frame spectra},
Ann. Fenn. Math. 46 (2021), 483--493; arXiv:2002.03855.
\url{https://arxiv.org/abs/2002.03855}.

\bibitem{LiLiu}
L. Li and B. Liu,
\emph{Fourier Frames on Salem Measures},
arXiv:2506.01280 (2025), especially the introduction and Section~10.
\url{https://arxiv.org/abs/2506.01280}.
\end{thebibliography}

\end{document}
